%% TECT Research Programme: Project Roadmap
%% Math01--Math36 Structured into Six Research Projects
%% Author: Jusang Lee | Date: April 2026

\documentclass[11pt,a4paper]{article}
\usepackage[margin=1in]{geometry}
\usepackage{amsmath,amssymb,amsthm}
\usepackage{graphicx}
\usepackage{bm}
\usepackage{hyperref}
\usepackage{braket}
\usepackage{array,booktabs}
\usepackage{xcolor}

\newcommand{\Nloop}{N_{\mathrm{loop}}}
\newcommand{\Dmu}{D_{\mu}}
\newcommand{\Gstar}{G_{*}}
\newcommand{\mstar}{m^{*}}
\newcommand{\HilL}{\mathcal{H}_{L}}

\newtheorem{theorem}{Theorem}[section]
\newtheorem{lemma}[theorem]{Lemma}
\newtheorem{openproblem}[theorem]{Open Problem}

\title{Topological Energy Condensate Theory (TECT):\\[4pt]
  \large Research Programme Roadmap --- Projects I through VI\\
  (Documents Math01--Math36)}
\author{Jusang Lee\\
  \small Independent Researcher, Incheon, Korea}
\date{April 2026}

\begin{document}
\maketitle
\tableofcontents
\newpage

% ---- Abstract ----
\section*{Abstract}
\addcontentsline{toc}{section}{Abstract}
Topological Energy Condensate Theory (TECT) derives gravity, gauge fields, and Standard
Model fermions from a single $\mathbb{C}^{3}$-valued energy condensate whose equilibrium
configuration is Body-Centered Cubic (BCC).
Documents Math01--Math36 decompose into six sequential projects:
(I) BCC ground-state theorem via $\mathbb{Z}_{2}$-symmetric Landau--Brazovskii theory
($\Nloop=6$ for BCC vs.\ $2$ for FCC and $0$ for SC);
(II) emergent $U(1)$ gauge structure;
(III) topological Weyl/Dirac fermion emergence;
(IV) microscopic Dirac carrier acceptance (Clifford closure, Gates 1--4);
(V) three-generation flavor structure and SM$+$GR infrared limit;
(VI) numerical PDE certification on the Class~II branch at
$(\mu^{2},\lambda,\gamma)=(0.26,-0.43,1.62)$.
\newpage

% =======================================================
\section{Overview: TECT Logical Chain and Project Map}
% =======================================================

The fundamental hypothesis is
\begin{equation}
\boxed{
  \text{Universe} = \text{BCC condensate } \Psi_{0}
  \;\xrightarrow{\text{fluctuations}}\;
  \text{gauge fields} + \text{fermions} + \text{gravity}.
}
\label{eq:master}
\end{equation}
The condensate field $\Psi(x)\in\mathbb{C}^{3}$ evolves under
\begin{equation}
\mathcal{L}_{\mathrm{mic}}
=\mathcal{L}_{\Psi}^{\mathrm{core}}[\Psi]
+\frac{M_{X}^{2}}{2}X_{i}^{a}X_{i}^{a}
+\alpha_{X}X_{i}^{a}J_{i}^{a}[\Psi]
+\beta_{X}X_{i}^{a}K_{i}^{a}[\Psi],
\label{eq:Lmic}
\end{equation}
with core operator
\begin{equation}
\mathcal{L}_{\Psi}^{\mathrm{core}}
=-\mu^{2}\Psi+Z\nabla^{2}\Psi-Y\nabla^{4}\Psi
-\lambda|\Psi|^{2}\Psi-\gamma|\Psi|^{4}\Psi.
\label{eq:Lcore}
\end{equation}

\subsection*{Project Status Summary}
\begin{center}
\renewcommand{\arraystretch}{1.45}
\begin{tabular}{clll}
\toprule
\textbf{Project} & \textbf{Files} & \textbf{Objective} & \textbf{Status} \\
\midrule
I   & Math01--05 & BCC ground-state uniqueness      & \textcolor{green!55!black}{Proved} \\
II  & Math06--09 & Emergent $U(1)$ gauge structure   & \textcolor{green!55!black}{Proved} \\
III & Math10--14 & Topological Weyl/Dirac fermions   & \textcolor{green!55!black}{Proved} \\
IV  & Math15--24 & Microscopic Dirac carrier audit   & \textcolor{orange!80!black}{Certificate pending} \\
V   & Math25--30 & Flavor structure \& SM+GR limit   & \textcolor{orange!80!black}{Framework complete} \\
VI  & Math31--36 & Numerical PDE certification       & \textcolor{red!70!black}{Active frontier} \\
\bottomrule
\end{tabular}
\end{center}

Logical dependency chain:
\begin{equation}
\text{I}\;\longrightarrow\;
\text{II}\;\longrightarrow\;
\text{III}\;\longrightarrow\;
\text{IV}\;\longrightarrow\;
\text{V}\;\longrightarrow\;
\text{VI}.
\end{equation}

% =======================================================
\section{Project I: BCC Ground-State Theorem (Math01--05)}
% =======================================================

\subsection{Core Objective}
Prove that the BCC Bravais lattice is the unique global minimum of the
$\mathbb{Z}_{2}$-symmetric Landau--Brazovskii free energy under pure quartic
(4-wave) fluctuations at 1-loop.

\subsection{Governing Free Energy}
\begin{align}
F[\{\Psi_{\mathbf{k}}\}]
&=\sum_{|\mathbf{k}|=k_{0}}\bigl[r+\kappa(k^{2}-k_{0}^{2})^{2}\bigr]|\Psi_{\mathbf{k}}|^{2}
\notag\\
&\quad+u\!\!\!\!\!\!\sum_{\substack{\mathbf{k}_{1}+\cdots+\mathbf{k}_{4}=0\\|\mathbf{k}_{i}|=k_{0}}}
\!\!\!\!\!\!\Psi_{\mathbf{k}_{1}}\Psi_{\mathbf{k}_{2}}\Psi_{\mathbf{k}_{3}}\Psi_{\mathbf{k}_{4}},
\label{eq:LB}
\end{align}
with $\mathbb{Z}_{2}$: $\Psi\to-\Psi$ forbidding cubic invariants.

\subsection{Main Theorem}
\begin{theorem}[BCC Selection --- Math01]
The non-trivial 4-wave resonance loop multiplicities are
\begin{equation}
\Nloop^{\mathrm{BCC}}=6,\quad
\Nloop^{\mathrm{FCC}}=2,\quad
\Nloop^{\mathrm{SC}}=0.
\label{eq:Nloop}
\end{equation}
At 1-loop, $\Delta F_{\mathrm{BCC}}<\Delta F_{\mathrm{FCC}}<\Delta F_{\mathrm{SC}}$,
and BCC is selected via a fluctuation-induced first-order transition.
\end{theorem}

\subsection{Robustness Extensions}
\begin{description}
\item[\textbf{Math02 (Universality A)}]
  Finite shell thickness $\Delta k/k_{0}\ll 1$ preserves
  $\Nloop^{\mathrm{BCC}}>\Nloop^{\mathrm{FCC}}$.

\item[\textbf{Math03 (Universality B)}]
  Amplitude variation $A_{i}=\bar{A}+\delta A_{i}$ shows equal-amplitude
  BCC is a strict local minimum.

\item[\textbf{Math04--05 (Multi-Shell)}]
  For $S=\bigcup_{a=1}^{K}S_{a}$, BCC selection extends under the
  single-shell dominance condition.
\end{description}

\subsection{Deliverables}
\begin{enumerate}
\item Exact $\Nloop$ for BCC, FCC, SC. \hfill\textcolor{green!55!black}{[Done]}
\item 1-loop ordering $\Delta F_{\mathrm{BCC}}<\Delta F_{\mathrm{FCC}}$. \hfill\textcolor{green!55!black}{[Done]}
\item Universality under $\Delta k$ and amplitude variation. \hfill\textcolor{green!55!black}{[Done]}
\item Multi-shell resonance network theorem. \hfill\textcolor{green!55!black}{[Done]}
\end{enumerate}

% =======================================================
\section{Project II: Emergent Gauge Structure (Math06--09)}
% =======================================================

\subsection{Core Objective}
Derive emergent $U(1)$ and higher gauge fields from low-energy fluctuations
around the BCC background $\Psi_{0}$.

\subsection{Transverse Decomposition}
Expanding $\Psi=\Psi_{0}+\delta\Psi$ and decomposing for each $\mathbf{k}$:
\begin{equation}
\delta\mathbf{u}=u_{L}\hat{\mathbf{k}}+u_{1}\mathbf{e}_{1}+u_{2}\mathbf{e}_{2}.
\end{equation}

\subsection{Main Theorems}
\begin{theorem}[$U(1)$ Emergence --- Math06]
In the IR limit $|\mathbf{k}|a\ll1$, the transverse sector carries
$SO(2)\simeq U(1)$. The emergent covariant derivative is
$\Dmu\sim\partial_{\mu}+iA_{\mu}$, $A_{\mu}\in\mathfrak{u}(1)$.
\end{theorem}

\begin{theorem}[$\mathbb{CP}^{2}$ Bundle --- Math07]
The orientation $z(x)\in\mathbb{C}^{3}$, $z^{\dagger}z=1$,
defines a principal $U(1)$ bundle over $\mathbb{CP}^{2}$.
The Berry connection yields $F=dA$.
\end{theorem}

\begin{theorem}[Dirac Coupling Non-Vanishing --- Math09]
For any $\Gstar\neq 0$, the reduced shell coefficients
$v_{\alpha}=e_{i}^{\alpha}(\mathsf{N}_{ij})^{a}(\Gstar)_{j}\neq0$, $\alpha=1,2$.
\end{theorem}

\subsection{Deliverables}
\begin{enumerate}
\item Emergent $U(1)$ from transverse BCC modes. \hfill\textcolor{green!55!black}{[Done]}
\item $\mathbb{CP}^{2}$ geometric bundle. \hfill\textcolor{green!55!black}{[Done]}
\item Non-vanishing Dirac coupling. \hfill\textcolor{green!55!black}{[Done]}
\item Full $SU(3)\times SU(2)\times U(1)$ from $\mathbb{C}^{3}$ bundle. \hfill\textcolor{orange!80!black}{[In progress]}
\end{enumerate}

% =======================================================
\section{Project III: Topological Fermion Emergence (Math10--14)}
% =======================================================

\subsection{Core Objective}
Establish massless Weyl/Dirac fermions at Bloch-band crossings on the
BCC reciprocal shell, protected by residual valley symmetry.

\subsection{Valley-Paired Block}
\begin{equation}
H_{D}(\mathbf{p})=
\begin{pmatrix}h_{+}(\mathbf{p})&0\\0&h_{-}(\mathbf{p})\end{pmatrix},
\quad h_{-}(\mathbf{p})=-h_{+}(-\mathbf{p}).
\end{equation}

\subsection{Main Theorems}
\begin{theorem}[Weyl Node --- Math10, Math11]
After shell-adapted coordinate change:
\begin{equation}
H_{D}(\mathbf{p})=\sum_{i=1}^{3}v_{i}p_{i}\alpha_{i},
\quad\alpha_{i}:=\tau_{z}\otimes\sigma_{i}.
\end{equation}
\end{theorem}

\begin{theorem}[Mass Protection --- Math12]
The residual $\mathbb{Z}_{2}^{\mathrm{valley}}$ forbids constant mass operators:
\begin{equation}
\det H_{D}(0)=0.
\end{equation}
\end{theorem}

\begin{theorem}[Little-Group Mismatch --- Math13]
Inequivalent representations of $H_{*}$ in $\mathcal{H}_{D}$ and $\mathcal{H}_{R}$
imply $V(0)=0$ (no remote-sector contamination at $p=0$).
\end{theorem}

\begin{theorem}[Clifford Closure --- Math14, Math18]
The full low-energy space
$\mathcal{H}_{\mathrm{low}}=\mathbb{C}^{2}_{\mathrm{int}}\otimes\mathbb{C}^{2}_{\mathrm{valley}}$
realises a minimal $4\times4$ Dirac representation.
\end{theorem}

\subsection{Deliverables}
\begin{enumerate}
\item Weyl node at BCC shell crossing. \hfill\textcolor{green!55!black}{[Done]}
\item Mass protection via $\mathbb{Z}_{2}^{\mathrm{valley}}$. \hfill\textcolor{green!55!black}{[Done]}
\item Little-group mismatch $\Rightarrow V(0)=0$. \hfill\textcolor{green!55!black}{[Done]}
\item Clifford closure on $\mathbb{C}^{4}$. \hfill\textcolor{green!55!black}{[Done]}
\end{enumerate}

% =======================================================
\section{Project IV: Microscopic Dirac Carrier Acceptance (Math15--24)}
% =======================================================

\subsection{Core Objective}
Finite operator audit confirming all four acceptance gates:
\begin{equation}
\boxed{
\text{Gate 1: isolation}
\;\wedge\;
\text{Gates 2--3: spectrum}
\;\wedge\;
\text{Gate 4: mass protection.}
}
\end{equation}

\subsection{Step-by-Step Programme}

\paragraph{Math15 (Audit Formulation).}
Open questions recast as finite operator audit on the linearized Bloch
operator around the locked branch. Central object: $P_{0}=z_{0}z_{0}^{\dagger}$.

\paragraph{Math16--17 (PDE Status).}
\begin{equation}
\boxed{\text{proof chain complete};\quad
\text{microscopic coefficient certificate: pending.}}
\end{equation}

\paragraph{Math18 (Clifford Closure).}
The reduced symbol on the doubled valley space yields the $4\times4$ Dirac
Hamiltonian. $\mathrm{Cl}(3,0)$ closes on $\{\tau_{z}\otimes\sigma_{i}\}_{i=1}^{3}\cup\{\tau_{x}\otimes\sigma_{0}\}$.

\paragraph{Math19--20 (Explicit Shell Basis).}
\begin{equation}
\mathcal{H}_{D}^{(N)}\cong
\mathrm{span}\bigl\{|\Gstar,\sigma,\tau\rangle\bigr\}_{\sigma,\tau\in\{+,-\}}.
\end{equation}

\paragraph{Math21 (Acceptance Milestone).}
Frontier identified as Gates 2--3 (spectrum audit) plus Gate 4 (mass protection).

\paragraph{Math22--23 (Non-Enriched Benchmark).}
$\det H_{\mathrm{non\text{-}enr}}^{(1)}(0)\neq0$ verified for $N=1$.

\paragraph{Math24 (Small-$p$ Drift Control).}
\begin{equation}
\|H^{(N)}(p)-H^{(N)}(0)\|\leq C_{\mathrm{drift}}\cdot|p|,\quad|p|\leq\rho.
\end{equation}

\subsection{Deliverables}
\begin{enumerate}
\item Gates 1--4 proved. \hfill\textcolor{green!55!black}{[Done]}
\item Explicit coefficient certificate. \hfill\textcolor{red!70!black}{[Pending --- Project VI]}
\end{enumerate}

% =======================================================
\section{Project V: Flavor Structure and SM+GR Infrared Limit (Math25--30)}
% =======================================================

\subsection{Core Objective}
Establish the three-generation left-handed space $\HilL$ and the
minimal theorem checklist for an SM$+$GR-type infrared description.

\subsection{Main Results}

\paragraph{Math25 (Local Uniqueness).}
The $3\times3$ Hessian of the reduced potential is positive-definite
at the selected branch: local uniqueness established.

\paragraph{Math26 (SM+GR Target Theorem).}
\begin{equation}
\boxed{\mathrm{TECT}_{\mathrm{IR}}\supseteq\mathrm{SM}+\mathrm{GR}.}
\end{equation}
Minimal checklist: 7 required lemmas for $SU(3)\times SU(2)\times U(1)$
minimally coupled to $g_{\mu\nu}$.

\paragraph{Math27 (Three-Generation Isolation, K1).}
\begin{equation}
\HilL=\mathrm{Ran}(P_{L}),\quad\dim\HilL=3.
\end{equation}

\paragraph{Math28 (Kinetic Gram Matrix, K2A).}
\begin{equation}
G_{IJ}=\langle L_{I}|\hat{K}|L_{J}\rangle,\quad I,J=1,2,3.
\end{equation}
Must be positive-definite for physical normalizability.

\paragraph{Math29 (BCC Ground State).}
\begin{equation}
\Psi_{0}(\mathbf{x})=\phi_{0}\!\!\sum_{\mathbf{G}\in\mathcal{S}_{\mathrm{BCC}}}
\!\!z_{\mathbf{G}}\,e^{i\mathbf{G}\cdot\mathbf{x}},\quad
|\mathcal{S}_{\mathrm{BCC}}|=12.
\end{equation}

\paragraph{Math30 (Completion Status).}
\begin{equation}
\boxed{\text{proof complete};\quad
\text{certificate incomplete (microscopic audit required).}}
\end{equation}

\subsection{Deliverables}
\begin{enumerate}
\item $\dim\HilL=3$ (three generations). \hfill\textcolor{green!55!black}{[Done]}
\item SM$+$GR checklist framework (7 lemmas). \hfill\textcolor{green!55!black}{[Done]}
\item Explicit $G_{IJ}$ entries. \hfill\textcolor{red!70!black}{[Pending --- Project VI]}
\end{enumerate}

% =======================================================
\section{Project VI: Numerical PDE Certification (Math31--36)}
% =======================================================

\subsection{Core Objective}
Complete the finite microscopic audit and certify the PDE coefficient
extraction on the Class~II working branch.

\subsection{Locked Parameters}
\begin{equation}
\boxed{(\mu^{2},\lambda,\gamma)=(0.26,\,-0.43,\,1.62).}
\label{eq:locked}
\end{equation}

\subsection{Sequential Steps}

\paragraph{Math31 (Enlarged Quartic Kernel).}
Four-class quartic kernel beyond the failed $(u,v,\kappa)$ truncation.

\paragraph{Math32 (Residual Quartic Closure).}
Exact four-class upgrade:
\begin{equation}
F_{4}=u_{1}\rho_{1}+u_{2}\rho_{2}+u_{3}\rho_{3}+u_{4}\rho_{4}.
\end{equation}

\paragraph{Math33--34 (PDE Coefficients + Seed Profile).}
\begin{equation}
\lambda_{\parallel}=\tfrac{\lambda_{c}}{2}I_{3},\quad
\alpha=\beta=-\frac{\beta_{X}\gamma_{\perp}}{M_{X}^{2}}I_{c}.
\end{equation}
Explicit seed profile inserted; Bloch matrix elements $\mathcal{U}_{j}(0)$ extracted.

\paragraph{Math35 (Canonical Longitudinal Coefficient).}
\begin{equation}
u_{\parallel}=\frac{1}{\phi_{0}}\,\hat{e}_{\parallel}^{\dagger}\,\mathcal{U}(0)\,\hat{e}_{\parallel}.
\end{equation}

\paragraph{Math36 (Class~II Elimination --- Active Frontier).}
Heavy mediator $X_{i}^{a}$ integrated out:
\begin{equation}
X_{i}^{a}=-\frac{1}{M_{X}^{2}}
\bigl(\alpha_{X}J_{i}^{a}+\beta_{X}K_{i}^{a}\bigr).
\label{eq:Xelim}
\end{equation}
The resulting effective Lagrangian is the starting point for full BCC linearization.

\subsection{Deliverables}
\begin{enumerate}
\item Four-class quartic kernel. \hfill\textcolor{orange!80!black}{[In progress]}
\item PDE coefficients $(\lambda_{\parallel},v_{\perp},v_{\parallel})$. \hfill\textcolor{orange!80!black}{[In progress]}
\item Positive-definite $G_{IJ}$ with explicit values. \hfill\textcolor{red!70!black}{[Target]}
\item Full BCC linearization after Class~II elimination. \hfill\textcolor{red!70!black}{[Target]}
\item $\mstar=0.3138$ derived analytically. \hfill\textcolor{red!70!black}{[Final target]}
\end{enumerate}

% =======================================================
\section{Immediate Next Step: Math37 Target}
% =======================================================

Five sequential steps to close Project~VI:
\begin{enumerate}
\item Insert Eq.~\eqref{eq:Xelim} into $\mathcal{L}_{\mathrm{mic}}$ to obtain
      the exact induced effective quartic Lagrangian.
\item Perform BCC Bloch decomposition on the 12-vector star
      $\mathcal{S}_{\mathrm{BCC}}$.
\item Extract $(v_{\parallel},v_{\perp},\lambda_{\parallel})$ numerically
      at parameters~\eqref{eq:locked}.
\item Verify $\det H_{D}(0)=0$ (Weyl node survives Class~II elimination).
\item Compute $\mstar$ analytically and verify $\mstar=0.3138\pm\delta\mstar$
      against numerical simulation.
\end{enumerate}

\begin{openproblem}[Analytical Derivation of $\mstar$]
Derive $\mstar=0.3138$ analytically from the BCC linearization of the
Class~II effective Lagrangian at $(\mu^{2},\lambda,\gamma)=(0.26,-0.43,1.62)$.
This constitutes the first analytical confirmation of a TECT-predicted
physical observable against numerical simulation.
\end{openproblem}

% =======================================================
\section{Overall Programme Status}
% =======================================================

\begin{align}
&\text{Projects I--III: all theorem-level proofs complete.}\notag\\
&\text{Project IV: Gates 1--4 proved; explicit coefficient values pending.}\notag\\
&\text{Project V: framework complete; explicit Gram matrix entries pending.}\notag\\
&\text{Project VI: active frontier. Math36 (Class~II elimination) is the immediate next step.}
\end{align}

\end{document}
